\section{The \texttt{IncoherentPhonon\_process} McStas Component}
A component to simulate inelastic scattering in the incoherent approximation
Takes into account one, two, and three phonon scattering explicitly,
and multi-phonon scattering (with n\textgreater{}4) via the saddle point method

\subsection*{Identification}
\begin{itemize}
  \item \textbf{Site:} 
  \item \textbf{Author:} Victor Laliena
  \item \textbf{Origin:} University of Zaragoza
  \item \textbf{Date:} 06.11.2018
\end{itemize}

\subsection*{Description}
\begin{lstlisting}
Part of the Union components, a set of components that work together and thus
expects geometry and physics within McStas.
The use of this component requires other components to be used.

1) One specifies a number of processes using process components like this one
2) These are gathered into material definitions using Union_make_material
3) Geometries are placed using Union_box / Union_cylinder, assigned a material
4) A Union_master component placed after all of the above


Algorithm:
Described elsewhere, see e.g. <a href="https://doi.org/10.3233/JNR-190117">https://doi.org/10.3233/JNR-190117</a>
\end{lstlisting}

\subsection*{Input parameters}
Parameters in \textbf{boldface} are required; the others are optional.

\begin{longtable}{p{0.22\textwidth}p{0.12\textwidth}p{0.46\textwidth}p{0.14\textwidth}}
\toprule
\textbf{Name} & \textbf{Unit} & \textbf{Description} & \textbf{Default} \\
\midrule
\endhead
T & K & Temperature & 394 \\
density & g/cm3 & Material density & 6.0 \\
M & amu & ion mass & 50.94 \\
sigmaCoh & barns & Coherent scattering cross section & 0.0184 \\
sigmaInc & barns & Incoherent scattering cross section & 5.08 \\
DOSfn & string & Path to the file that contains the DoS & NULL \\
nphe\_exact & 1 & Number of terms in the phonon expansion taken exact. Has to be between 1 and 3. & 1 \\
nphe\_approx & 1 & Number of terms in the phonon expansion taken approximate. & 0 \\
approx & 1 & Approximation type: 0 gaussian, 1 saddle point & 0 \\
mph\_resum & 0/1 & Resumate the remaining terms of the phonon expansion via a saddle point: 0 No, 1 Yes & 0 \\
nxs & 1 & Number of energy points at which the total cross sections are precomputed & 1000 \\
kabsmin & A\textasciicircum{}-1 & Lower cut-off for the neutron wave-vector k & 0.1 \\
kabsmax & A\textasciicircum{}-1 & Higher cut-off for the neutron wave-vector k & 25 \\
interact\_fraction & 1 & How large a part of the scattering events should use this process 0-1 (sum of all processes in material = 1) & -1 \\
init & string & Name of Union\_init component (typically "init", default) & "init" \\
\bottomrule
\end{longtable}

\subsection*{Links}
\begin{itemize}
  \item \href{run:/home/willend/willend-McCode/mcstas-comps/union/IncoherentPhonon_process.comp}{Source code} for \texttt{IncoherentPhonon\_process.comp}.
  \item See \textless{}a href="https://doi.org/10.3233/JNR-190117"\textgreater{}https://doi.org/10.3233/JNR-190117\textless{}/a\textgreater{}
\end{itemize}
\IfFileExists{IncoherentPhonon_process_static.tex}{\input{IncoherentPhonon_process_static.tex}}{}